% HD-S46. GPT-6 Astra Pro /root/sixfolds, inherited effort; 21 November 2026.
\section{HD-S46: intrinsic three-dimensional blocks}
\label{sec:three-block-kernel-scope}

Let \(K\subset\mathbf C\) have characteristic zero. Let
\(N\subset GL_{3r}(K)\), \(r\ge5\), be finite and absolutely
irreducible, generated by actual matrices \(x\) satisfying
\(\operatorname{rank}(x-I)\le3\). Suppose that its complex
representation has a system of \(r\) blocks of dimension three.
Then the kernel \(T\) of its action on these blocks is
characteristic in \(N\). Each of the original block projections
has entries in \(K\). Every actual \(K\)-linear group normalizing
\(N\) therefore permutes these same descended three-dimensional
\(K\)-blocks.

For the natural sixfold sources under consideration, take
\(N=N_3\), the normal subgroup generated by all actual
codimension-three elements of the coefficient group \(G\).
In the absolutely irreducible branch of rank at least fifteen,
the three-dimensional complex block alternative consequently
gives an actual rank-three quotient source. At period rank nine
its ordinary Hermitian determinant is among the already supplied
classes: \([-1]\) or \([-3]\) for \(K=\mathbf Q(i)\), and
\([-1]\) for
\[
 K=\mathbf Q(\sqrt{-d}),\qquad
 d\in\{2,7,11,19,43,67,163\}.
 \tag{S46.1}
\]
The structural descent statement applies to every characteristic-zero
\(K\). The displayed determinant conclusion has exactly the
eight-field scope of HD-S40 and HD-G27. The one-dimensional and
two-dimensional complex block alternatives are not asserted to
be settled by this argument.

\paragraph{1. Actual swaps and a simultaneous tree gauge.}
Write \(V_{\mathbf C}=\bigoplus_{i=1}^rV_i\), with
\(\dim V_i=3\). If an element cycles \(\ell\) blocks,
a fixed vector on their sum is determined by its value in one
block and must be fixed by the product of the transition maps.
The cycle contributes at least \(3(\ell-1)\) to its fixed-space
codimension. Thus every generator outside \(T\) exchanges
exactly two blocks, is the identity on all the others, and is
literally
\[
 x|_{V_i\oplus V_j}=
   \begin{pmatrix}0&A\\ A^{-1}&0\end{pmatrix}.
 \tag{S46.2}
\]
Indeed its entire codimension budget is already three; the
return map must fix all three coordinates. In particular this
generator is an involution. Every generator inside \(T\)
has nonidentity coordinates on at most three blocks.

The graph of the exchanged pairs of the outside generators is
connected: a disconnected component sum would be invariant under
all the generators. Choose a spanning tree and one actual generator
on each edge. Choose \(b_1:\mathbf C^3\longrightarrow V_1\).
Along an oriented edge from an already based block \(i\) to a new
block \(j\), put \(b_j=x_{ij}|_{V_i}b_i\). The inverse relation
in (S46.2) makes this tree generator swap the two bases literally.
Under the specified linear isomorphism
\[
 B:(\mathbf C^3)^r\longrightarrow V_{\mathbf C},
       (u_i)\longmapsto\sum_i b_i(u_i),
 \tag{S46.3}
\]
the chosen generators are pure transpositions. Tree transpositions
generate all of \(S_r\): conjugating along a tree path gives the
transposition of its endpoints. Hence the block permutation image
is exactly \(S_r\), and in these actual coordinates
\[
                         N=T\rtimes S_r.
 \tag{S46.4}
\]
No assertion that the coordinate map \(B\) is defined over \(K\)
has been used.

\paragraph{2. Three-support generators and localized commutators.}
For an outside generator whose permutation is \((i\ j)\),
multiply it by the inverse of the pure transposition in (S46.4).
The result is an element of \(T\) supported on those two blocks.
Keep the inside generators unchanged. Let \(\mathcal D\) be
this set. The subgroup generated by its conjugates under the
pure \(S_r\) is normalized by \(S_r\), and together with
\(S_r\) contains all the original generators. Its permutation
kernel is therefore all of \(T\). In particular
\[
 T=\langle p d p^{-1}:p\in S_r,\ d\in\mathcal D\rangle,
 \qquad \text{each displayed generator has support at most three}.
 \tag{S46.5}
\]
Let \(L\subset GL_3(\mathbf C)\) be the image of \(T\)
on one block. Pure permutations identify all the images with
this same finite group, so \(T\subset L^r\). The group \(L\)
is irreducible: a proper common invariant subspace repeated in
every block would be invariant under \(T\) and the pure
permutations, contradicting irreducibility of \(N\).

We have the literal subgroup inclusion
\[
                              (L')^r\subset T\subset L^r.
 \tag{S46.6}
\]
To prove it, take two coordinate entries \(a,b\) of the
generators in (S46.5). Place their chosen coordinates in block
one by pure permutations. Each has at most two other support
blocks. Since \(r\ge5\), further permutations fixing block one
place these other support sets disjointly. The commutator is
then exactly \(\operatorname{diag}([a,b],I_3,\ldots,I_3)\).
The subgroup of first entries occurring in elements supported
only on block one is normal in \(L\): conjugate by any element
of \(T\) with a prescribed first entry. The entries of the
generators in (S46.5) generate \(L\), so the normal closure
of their pairwise commutators is \(L'\). This proves the first
factor of (S46.6); the pure permutations give every other factor,
and disjointly supported factors commute as matrices.

\paragraph{3. The intrinsic description of the kernel.}
The only new classification input is the finite primitive
\(PGL_3(\mathbf C)\) list of Dolgachev--Iskovskikh,
\emph{Finite subgroups of the plane Cremona group}, Theorem4.8.
The chapter is published in \emph{Algebra, Arithmetic, and
Geometry}, Progress in Mathematics269 (2009),443--548,
DOI \url{https://doi.org/10.1007/978-0-8176-4745-2_11}.
The author-hosted typeset proof was read at its printed p.459,
PDF page23, at
\url{https://math.nyu.edu/~tschinke/.manin/final/dolgachev/dolgachev.pdf}.
Its pagination starts at437 rather than the published443; we
retain its literal page locator rather than identifying those
pagination systems. The theorem lists \(A_5,PSL_2(7),A_6\)
and the solvable Hessian groups
\(3^2:SL_2(3),3^2:C_4,3^2:Q_8\).
The binary tetrahedral group \(SL_2(3)\) has normal \(Q_8\)
with quotient \(C_3\), so all the stated Hessian groups are
solvable.

If \(L\) is imprimitive, its irreducible three-dimensional
representation is monomial on three lines. The diagonal kernel
is abelian and the permutation quotient lies in \(S_3\), so
\(L\) is solvable. In the primitive case the cited list, applied
to its literal projective image, shows that either \(L\) is
solvable or
\[
             L/Z\simeq S\in\{A_5,PSL_2(7),A_6\},
             \qquad Z=L\cap\mathbf C^*I_3.
 \tag{S46.7}
\]
The scalar kernel is abelian, so these alternatives exhaust \(L\).
The three displayed groups \(S\) are nonabelian simple and perfect.

Write \(\operatorname{Rad}(F)\) for the largest solvable normal
subgroup of a finite group. It is characteristic: products of
solvable normal subgroups are solvable, so their finite product
is the unique largest one. Also \(\operatorname{Rad}(S_r)=1\)
for \(r\ge5\), since its nontrivial normal subgroups contain
the nonabelian simple \(A_r\).

When \(L\) is solvable, (S46.6) makes \(T\) solvable,
and every solvable normal subgroup of \(N\) maps trivially
to \(N/T=S_r\). Consequently
\[
                         T=\operatorname{Rad}(N).
 \tag{S46.8}
\]
When (S46.7) holds, the image of \(L'\) in \(S\) is
all of \(S\). Projectivizing each coordinate therefore gives,
by (S46.6), the exact surjection
\[
 \pi:T\longrightarrow S^r,
       \qquad\ker\pi=Z_0:=T\cap Z^r.
 \tag{S46.9}
\]
The kernel \(Z_0\) is abelian and normal in \(N\). Every
solvable normal subgroup of \(N\) lies in \(T\); its image
is a solvable normal subgroup of \(S^r\), hence trivial.
For this last assertion, normal subgroups of a product of
nonabelian simple groups are products of its factors: a nontrivial
projection to a factor, commuted with that factor, gives the
entire factor inside the normal subgroup, and one continues
with the other projections. Thus
\[
 \operatorname{Rad}(N)=Z_0,\qquad
 N/Z_0\xrightarrow{\sim} S^r\rtimes S_r,
        \quad(tp)Z_0\longmapsto(\pi(t),p).
 \tag{S46.10}
\]
This is an actual isomorphism: (S46.4) gives uniqueness of
\(tp\), pure conjugation permutes the coordinates of \(\pi(t)\),
and (S46.9) gives surjectivity and the exact kernel.

The base \(S^r\) is the unique minimal nontrivial normal
subgroup of this quotient. It is minimal because its normal
subgroups are products of factors and the top \(S_r\) is
transitive. Its centralizer is trivial: a centralizing element
has trivial permutation component and each remaining coordinate
is central in \(S\), hence trivial. Any other minimal normal
subgroup would have trivial intersection with the base, and
would commute with it, a contradiction. It follows that in
the nonsolvable case \(T\) is the inverse image in \(N\)
of the unique minimal normal subgroup of
\(N/\operatorname{Rad}(N)\). This and (S46.8) prove that
\(T\) is characteristic in every case, including \(r=5,6\).

\paragraph{4. Actual descent of each block projection.}
Each \(V_i\) is an irreducible \(T\)-module. The group \(L\)
is not abelian, since finite commuting matrices have common
eigenlines; thus \(L'\ne1\). Choose \(a\in L'\setminus\{1\}\).
By (S46.6), \(T\) contains the element acting by \(a\)
on \(V_i\) and by the identity on all other blocks. This
distinguishes \(V_i\) from every other \(V_j\) as a
\(T\)-module. The \(r\) block modules are consequently
pairwise inequivalent and occur with multiplicity one.

Let \(\chi_i\) be the character on \(V_i\). In the original
matrix coordinates, define
\[
       p_i=\frac{3}{|T|}\sum_{t\in T}\chi_i(t^{-1})\rho(t).
 \tag{S46.11}
\]
The character is a class function, so \(p_i\) is scalar on
each irreducible \(V_j\). Averaging the action on
\(\operatorname{Hom}(V_i,V_j)\) gives
\(|T|^{-1}\sum_t\chi_j(t)\chi_i(t^{-1})=\delta_{ij}\).
The trace of \(p_i|_{V_j}\) is therefore \(3\delta_{ij}\),
so it is exactly the projection onto \(V_i\). In particular
\(p_ip_j=\delta_{ij}p_i\), \(\sum_i p_i=I\), and
\(\operatorname{rank}p_i=3\).

Let \(E=K(\zeta_h)\), where \(h\) is the exponent of \(T\).
All character values in (S46.11) belong to \(E\). The centre
of the actual image algebra \(E\rho(T)\) is exactly
\(\bigoplus_i Ep_i\): an element in the centre acts by a
scalar on each irreducible block, and the displayed projectors
realize all such scalars. These are precisely its primitive
central idempotents. The finite Galois group of \(E/K\)
therefore permutes the \(p_i\). Its permutation commutes with
the natural \(S_r\)-action of \(N\), because every matrix
\(\rho(n)\) is defined over \(K\). The centralizer of
\(S_r\) on its \(r\) labels is trivial: the stabilizer
of label \(i\) fixes only \(i\), so a commuting permutation
must fix \(i\). It follows that
\[
       p_i\in M_{3r}(K),\qquad U_i=p_iK^{3r},\qquad
       U_i\otimes_K\mathbf C=V_i.
 \tag{S46.12}
\]
Every actual \(K\)-linear normalizer of \(N\) normalizes
the characteristic subgroup \(T\), hence permutes these same
primitive central projectors. This proves the structural assertion.

\paragraph{5. The exact coefficient and polarized quotient maps.}
Now let \(K\) be imaginary quadratic, let \(G\) be the
original finite coefficient group, with exact character field
\(K\), and let \(h\) be its prescribed positive invariant
Hermitian form, conjugate-linear in the first argument.
The inequivalent \(T\)-blocks are orthogonal: a cross term
would induce an intertwiner between distinct irreducibles.
Choose \(U=U_1\), let \(H\) be its stabilizer, and let
\(\sigma:H\to GL_K(U)\) be the restricted action. It is
absolutely irreducible because its \(T\)-action already is.
The induced-representation map is the actual map
\(t\otimes u\mapsto\rho(t)u\). Its pushforward-function
version and inverse are
\[
\begin{aligned}
 F_v(x)&=p_1\rho(x)v,& F_v(hx)&=\sigma(h)F_v(x),\\
 F_{\rho(g)v}(x)&=F_v(xg),&
 F&\longmapsto\sum_{H\backslash G}\rho(x)^{-1}F(x).
\end{aligned}
 \tag{S46.13}
\]
The inversion label \(Ht^{-1}\leftrightarrow tH\) relates
left-deck sheets to blocks. These inverse maps commute with
local inertia invariants: along a cycle of sheets a fixed
section is determined by one vector fixed by the product of
the transition maps. Thus the isomorphism extends to the
middle extensions and their parabolic cohomology at every
ramified stalk. The induced trace formula gives
\(\mathbf Q(\operatorname{tr}\rho)\subset
\mathbf Q(\operatorname{tr}\sigma)\subset K\), so the
faithful image of \(\sigma\) has exact field \(K\).

Put \(H_0=\ker\sigma\). From the original connected
\(G\)-cover form the actual curves and map
\[
 Y=\widetilde C/H,\qquad C'=\widetilde C/H_0,
       \qquad\alpha:\widetilde C\longrightarrow C'.
 \tag{S46.14}
\]
For \(v\in U(K)\setminus\{0\}\), put \(q=h(v,v)\).
Orthogonality and \(|G|=r|H|\) give equality of the actual
rational corner idempotents in \(\mathbf Q[G]\):
\[
 \frac{2(3r)}{|G|q}\sum_{g\in G}
       \operatorname{Re}h(v,\rho(g^{-1})v)g
 =\frac{2\cdot3}{|H|q}\sum_{h_0\in H}
       \operatorname{Re}h(v,\sigma(h_0^{-1})v)h_0.
 \tag{S46.15}
\]
Every coefficient off \(H\) is zero. Averaging over the
kernel \(H_0\) identifies the right side with the same
corner of \(\mathbf Q[H/H_0]\). An integer multiple
therefore defines the same abelian image inside
\(J\widetilde C\), and pullback identifies the primitive
factor of \(JC'\) with that image up to the actual isogeny.
For the left deck actions we use their pushforward homomorphisms
on the Jacobians, exactly as in HD-S23 and HD-G18. The identities
for the quotient morphism are
\[
 \alpha_*\alpha^*=[|H_0|],\qquad
       (\alpha^*)^*\theta_{J\widetilde C}=|H_0|\theta_{JC'}.
 \tag{S46.16}
\]
The real coefficient pairings have equal Schur constant
\(2|G|/((3r)q)=2|H|/(3q)\), so this is an identification
of the specified paired Hodge factors, retaining the corner
field action.

For completeness, the initially prescribed primitive line
may be \(Kw\) with \(w\notin U\). Put \(q_w=h(w,w)\)
and \(A(x)=v h(w,x)/q_w\). The matrix-coefficient element
\[
 b=\frac{2(3r)}{|G|}\sum_{g\in G}
          \operatorname{Re}\operatorname{tr}
                   (A\rho(g^{-1}))g
 \tag{S46.17}
\]
acts by \(A\) and its conjugate on the two coefficient
blocks, and by zero on the other rational blocks; this follows
from the same matrix-coefficient averaging that proves
(S46.11). Its Rosati identity is
\(b^\dagger b=(q/q_w)e_w\). Clearing its rational
coefficients by an integer \(M>0\) gives an actual isogeny
between the two primitive images with pullback polarization
multiplier \(M^2q/q_w\). Thus (S46.14)--(S46.17) apply
to the original specified factor, not only to a newly chosen
line. A positive rational multiplier changes its rank-six
Hermitian determinant by its sixth power, which is the norm
of its cube in \(K\).

\paragraph{6. The rank-three period endpoint.}
The faithful coefficient system on \(Y\) has rank three
and first parabolic cohomology of dimension six over \(K\).
Let \(g=g(Y)\), and let \(b\) count its genuine nontrivial
coefficient branch values, with codimensions \(c_j\ge1\).
Absolute irreducibility kills both invariant spaces, and the
exact Euler formula gives
\[
                       6=3(2g-2)+\sum_{j=1}^b c_j.
 \tag{S46.18}
\]
For \(g=1\), this gives \(b\le6\) and a pointed-base
parameter dimension at most six. For \(g=2\), it gives
\(b=0\) and dimension three; larger genus is impossible.
For \(g=0\), it gives \(b\le12\), with pointed-base
dimension \(b-3\). Fixing the finite monodromy labels leaves
a finite discrete set of covers over each marked base. After
a finite local parameter cover the actual quotient family
therefore factors through these pointed-curve moduli. The
isogenies (S46.16)--(S46.17) preserve the period rank.

Consequently period rank nine forces \(g=0\), \(b=12\),
and every \(c_j=1\): the rank-three coefficient group is
actually generated by twelve reflections. HD-S31/HD-S40 for
the Gaussian field and HD-G27 for the seven fields in
(S46.1) apply to this literal quotient group, giving exactly
the asserted supplied determinant classes. No equality of
dimensions has been used to infer dominance. This proves the
three-dimensional block reduction through the original
coefficient, curve, corner and polarization maps.
